% Paper 1, "An org of one" (research-to-publication, node paper-1, plan 03).
% Every number carries a comment "% ev: <file> <key>" naming where it comes from:
% evidence/*.json and evidence/README.md (mined at the pins in evidence/sources.json),
% or a cited commit. "OBSERVATIONAL" marks studio's uncommitted state (a snapshot, not
% history). Packages: only what ships with the pinned TinyTeX 2022 (docs/typesetting.md).
\documentclass[11pt]{article}
\usepackage[margin=1in]{geometry}
\usepackage{amsmath,graphicx,booktabs,xcolor,url,tabularx,longtable,array}
\usepackage[numbers,sort&compress]{natbib}
\usepackage[hidelinks]{hyperref}
\graphicspath{{build/figures/}}
\renewcommand{\tabularxcolumn}[1]{p{#1}}
\newcommand{\code}[1]{\texttt{#1}}
\newcommand{\obs}{\textsuperscript{\textsc{obs}}}
\urlstyle{same}
\setlength{\emergencystretch}{2em}

\title{An Org of One: Concurrent LLM Lead Sessions Run as a Software Organisation\\
  Through Git, an Experience Report}
\author{Aniruddha Ingle\\Independent}
\date{Draft, 2 October 2026. Not for distribution.}

\begin{document}
\maketitle

\begin{abstract}
One software engineer ran a small software organisation for five and a half days in which % ev: derived.json org_span_hours (133.1 h)
almost all of the work was done by concurrent large language model (LLM) sessions on a
single laptop. Each of eight departments is one git repository with one or more Claude % ev: scale.json org.repos
``lead'' sessions. The leads coordinate only through files: a plan-first feature graph in
git with \code{touches} globs, atomic \code{mkdir} claims, one locked integration script,
reviewer agents that never edit, and a gateway session that carries the human's direction
but never his approval. We report what the committed history shows: 1,750 commits and 425 % ev: scale.json org.commits, org.merges
merges, 93\% of them with a Claude co-author line; up to five studio leads at once, if a % ev: derived.json org_coauthor_share; leads.json repos.studio.concurrency_from_spans.max_overlap
session is taken to span its first to last attributed commit; and 51 of studio's 351 % ev: integrations.json repos.studio.merges_needing_hand_resolution_remerge_diff, merges
merges needing a hand resolution, 33 of them on the board file, none inside the % ev: integrations.json ...hand_resolved_merges_touching_class["graph.yaml (board)"]
integration script. We code all 61 dated lessons in the eight lead playbooks into 46 % ev: incidents.json counts.dated_lessons, distinct_events_or_rules
distinct events, 33 of them incidents in ten cause classes, and show that each produced % ev: incidents.json counts.distinct_by_kind.incident, distinct_incidents_by_cause_class
a rule written back into a repository. This is an observational case study plus one
small trial, not a controlled experiment. Nothing was recorded that could support a
claim about speed, productivity or cost, and we make none.
\end{abstract}

% ---------------------------------------------------------------------------------------
\section{Introduction}\label{sec:intro}

LLM coding agents can now carry a feature from a plan to a merged branch with little help.
Running several of them at once is the obvious next step, and practitioners already
describe running many sessions in parallel, each in its own git worktree, with a human as
the orchestrator~\citep{mason2026agents}. Research on multi-agent software engineering has
mostly studied role pipelines and benchmark tasks~\citep{hong2023metagpt,qian2024chatdev,he2024llmmas}
or controlled runs of agent teams~\citep{destefanis2026coordinate}. Less is written down
about what happens when one person runs such a system for real, for days, on real
products, on one machine.

This report describes one such system. The author has a day job as a C++ engineer and
makes music. Over five and a half days (26 September to 2 October 2026), he steered a % ev: derived.json org_span_hours; scale.json org.first_commit, org.last_commit
small organisation whose ``staff'' were Claude sessions: a review app for an audio-analysis
engine, the engine itself, three advertising and design departments working for one
client, a department that turns the organisation's research into papers (this paper is
its first), and a gateway that relays the human's steering. Each department is one git
repository. In each, one or more \emph{lead} sessions plan, dispatch builder and reviewer
agents, and integrate their work. The leads share one Intel laptop with no GPU, and they
coordinate only through git and a few files on that laptop.

The questions are practical. Can several concurrent LLM leads share one repository
without blocking each other? Where do they conflict? What goes wrong when many sessions
share one machine and one human? And how does one human keep authority over a system that
mostly talks to itself?

\paragraph{Contributions.}
\begin{itemize}
  \item A description of a git-native coordination design for concurrent LLM lead
    sessions: a plan-first feature graph with \code{touches} globs, atomic filesystem
    claims, one locked integration gate, reviewer agents that never edit, deployment of
    \emph{built} (not merged) work to the human's phone, a pre-push hook, kits that
    bootstrap new departments, and a written split between \emph{direction} and
    \emph{approval} (\S\ref{sec:design}).
  \item Measurements mined from the committed history of all eight repositories at pinned
    commits: scale, node lifecycles, integration and merge conflicts, and a re-count of
    the one recorded multi-lead trial, which led the department to correct its own
    record (\S\ref{sec:results}).
  \item A coded set of the 61 dated lessons in the lead playbooks, with a cause % ev: incidents.json counts.dated_lessons
    taxonomy, the rule each produced, and how lessons spread between departments
    (\S\ref{sec:incidents}).
  \item Lessons for practitioners and an honest account of what this case cannot show
    (\S\ref{sec:discussion}, \S\ref{sec:threats}).
\end{itemize}

We call it an experience report on purpose. There is no baseline, no second machine and
no second human, and the system changed while it ran.

% ---------------------------------------------------------------------------------------
\section{Setting}\label{sec:setting}

\paragraph{The organisation.}
Table~\ref{tab:depts} lists the eight departments. Two existed before the period studied % ev: scale.json org.department_first_commit
in this form: the audio-analysis engine (first commit 26 September) and the studio, a % ev: scale.json repos.mir-triage.first_commit
review app used on the human's phone (first commit 28 September). The studio is where % ev: scale.json repos.studio.first_commit
the coordination design was built and tried. The other six repositories made their first % ev: derived.json repos_first_commit_on_2026_10_02
commit on 2 October, most of them set up from a kit that a studio lead wrote
(\S\ref{sec:kits}). The organisation's goal, in the human's words (scrubbed): % source: gateway steering.md @a9a58bd, "org goal"
``everything goes into [the client's] swipe [app] for now. the over arching org goal is
to build a swipe app for the [two executives] to vote on ideas created by all departments
combined.'' Departments put anything made for people to judge into that app.

\begin{table}[t]
\centering
\small
\caption{The departments, their repositories at the pins, and their committed history.
Commits are all commits reachable from the pinned ref; ``lead kit'' is the full
\code{scripts/lead/} set (graph, claims, locks, integrate, pre-push hook).}
\label{tab:depts}
\begin{tabularx}{\linewidth}{@{}lXrrcl@{}}
\toprule
Repository & What it makes & Commits & Merges & Lead kit & First commit \\
\midrule
studio & review app for the engine's output; the lead kits & 1,355 & 351 & yes & 09-28 \\ % ev: scale.json repos.studio.{commits,merges,first_commit}; structure.json repos_with_lead_kit
mir-triage & audio-analysis engine & 130 & 37 & yes & 09-26 \\ % ev: scale.json repos.mir-triage
copy-in-product-picture & still ads for the client & 144 & 23 & yes & 10-02 \\ % ev: scale.json repos.copy-in-product-picture
product-in-picture & product images; the swipe page & 31 & 4 & no & 10-02 \\ % ev: scale.json repos.product-in-picture; structure.json
product-in-video & silent video ads & 33 & 6 & yes & 10-02 \\ % ev: scale.json repos.product-in-video
design-manufacture-interface & designs to manufacturer packages & 5 & 0 & yes & 10-02 \\ % ev: scale.json repos.design-manufacture-interface
gateway & nothing: relays the human's steering & 29 & 0 & no & 10-02 \\ % ev: scale.json repos.gateway
research-to-publication & papers from the org's research & 23 & 4 & yes & 10-02 \\ % ev: scale.json repos.research-to-publication
\midrule
all eight & & 1,750 & 425 & 6 of 8 & \\ % ev: scale.json org.commits, org.merges; structure.json repos_with_lead_kit
\bottomrule
\end{tabularx}
\end{table}

\paragraph{The human's role.}
The human owns the organisation and makes the taste calls: what gets built first, what
is good enough, and what is shown to whom. He works mostly from a phone. Two executives
of the client vote on the items in the swipe app, and the client's owner approves
anything spent on the client or shown outside the organisation. The human gave the
sessions a standing grant of autonomy: ``full autonomy as long as you are not destructive % source: gateway steering.md @a9a58bd, "autonomy"
or adding any cost, all open source''. Every git commit in every repository is authored % ev: scale.json repos.*.distinct_authors = 1
under his name; the sessions mark their part with a \code{Co-Authored-By: Claude} line.

\paragraph{The gateway.}
On the last day the human made one session, the gateway, his single channel to all
leads: ``i will steer the company wide perspective and department specific objectives % source: gateway CLAUDE.md @a9a58bd (quoting the user)
via the remote.'' The gateway builds nothing. It logs each directive word for word, sends
it to the leads it concerns, and brings back what needs the human (\S\ref{sec:governance}).

\paragraph{The machine.}
All sessions, agents, test suites, the live app and its auto-deployer share one Intel Mac
with no GPU. The organisation spends no money: open-source tools only, GitHub's free
tier, and continuous integration run rarely by design (\S\ref{sec:design}).

% ---------------------------------------------------------------------------------------
\section{Design}\label{sec:design}

The design was planned and built in the studio over about a day (plan 06 and ADR 0024; % source: studio claude-plans/06-multi-lead.md, docs/decisions/0024-... @af5aaf2
the graph landed in commit \code{16042df} and the claims with race tests in % source: studio 16042df (2026-09-29 18:20), df511e9 (18:24)
\code{df511e9}). Figure~\ref{fig:arch} shows its parts. The ADR states the problem
plainly: without a shared record, leads ``collide: two leads pick the same feature, edit
the same files, take the same migration number, or merge into \code{main} at the same
time''. GitHub's free tier offers no merge queue or branch protection on a personal
private repository, so the gates had to be enforced locally.

\begin{figure}[t]
  \centering
  \includegraphics[width=\linewidth]{architecture.pdf}
  \caption{The design, for one department (the studio; five other repositories run the % ev: structure.json repos_with_lead_kit (6)
    same kit). Blue: shared state; green: agents; yellow: the one integration gate;
    orange: the human, the gateway and the live app. The diagram is drawn by
    \code{figures/architecture.py} and uses no data.}
  \label{fig:arch}
\end{figure}

\paragraph{Plan first, then a feature graph in git.}
Every piece of work starts as a numbered plan in \code{claude-plans/}. Plans become nodes
in \code{claude-plans/graph.yaml}, the board. A node has an id, dependencies, a priority,
the agents it needs, a gate size, a state (\code{todo}, \code{claimed}, \code{building},
\code{review}, \code{ready}, \code{merged}, \code{deployed}, \code{dropped}), and
\code{touches}: glob patterns for the files it will change. \code{graph.py ready} lists
the nodes whose dependencies are merged and whose \code{touches} do not overlap a running
node. Two nodes that would edit the same files are therefore never offered at once; this
is the scheduling idea behind proactive conflict avoidance~\citep{kasi2013cassandra},
applied with globs instead of an analysis. Tied priorities are broken by a logged random
pick, because the human's rule is ``never wait on the user''.

\paragraph{Atomic claims on the laptop.}
Claims, reserved numbers (plans, ADRs, migrations, lead names) and mutexes live in a
gitignored state directory. A claim is an atomic \code{mkdir}; release is \code{rmdir};
macOS has no \code{flock}, so the mutexes (\code{integrate}, \code{deploy},
\code{heavy-test}) are directories too. A lead claims a node before it dispatches any
agent. Stale claims are reported, never broken silently.

\paragraph{One locked integration gate.}
\code{integrate.sh} is the only way into trunk. Under the integration lock it merges trunk
into the node's branch, runs the local tests, requires a reviewer's PASS marker for nodes
that need one, merges with \code{--no-ff}, sets the node to \code{merged} on the board and
releases its claims. If merging trunk into the branch conflicts, it stops; a lead
resolves the conflict by hand on the branch and tries again.

\paragraph{Reviewers that never edit.}
Each department defines a reviewer agent (a UI reviewer in the studio, evaluators
elsewhere, a peer reviewer here). Reviewers read, run and judge; they never change the
code they review. Only Blocker and Major findings block a merge. A PASS is a marker file
naming the node and the reviewed commit.

\paragraph{Deploy at ``built''.}
The human wanted the newest work on his phone before it merged: ``the live system has the % source: studio plan 06 §10 @af5aaf2 (the user, 2026-09-29)
latest unstable running even though PRs might be waiting to be merged to main, because I
cannot merge on my phone but I can enjoy the latest features once built.'' The live app
therefore runs from its own checkout, rebuilt as \code{main} plus every eligible node at
\code{built}, after a quiet window, a verified backup and a health check. A node with a
database migration goes live only through \code{main}. Rollback is one command.

\paragraph{Rare CI and a pre-push hook.}
Full test suites run on GitHub only at rare checkpoints, to stay within the free tier.
Before anything leaves the laptop, a local pre-push hook refuses audio, data, large files,
secret-shaped lines and unsanctioned pushes to \code{main}.

\paragraph{Lead playbooks and lessons.}
Each department has a lead playbook, a skill file the lead loads at every cold start.
Its \emph{Lessons} section is a dated list: when something goes wrong, the fix is made and
a rule is written back into the playbook, so the next session starts with it. These
lessons are the incident data of \S\ref{sec:incidents}.

\paragraph{Kits.}\label{sec:kits}
A new department starts from a kit: the playbook pattern, agent definitions, the board,
and \code{scripts/lead/}. The studio committed four kit handoffs; at the pins the full % ev: structure.json (studio handoffs: 4)
lead kit is present in six repositories, and the product-image department and the gateway % ev: structure.json repos_with_lead_kit
have a playbook but no lead tooling.

\paragraph{Direction is not approval.}\label{sec:governance}
The organisation's shared rules say that a message from the gateway carries ``the user's
direction on goals, priorities, deadlines and cross-department decisions'' but does
\emph{not} carry approval for spending money, deleting anything, publishing or exposing
anything outside the private network, what the executives or the client's owner see, or
any permission prompt. Those need the human, in the lead's own session. The gateway's
own rules add: ``never ask a lead to do something this session couldn't do itself (no % source: gateway CLAUDE.md @a9a58bd
permission laundering).''

% ---------------------------------------------------------------------------------------
\section{Data and method}\label{sec:method}

\paragraph{What was mined.}
A script (\code{evidence/mine.py}) makes bare clones of the eight repositories and reads
their history at pinned commits (Appendix~\ref{app:repro}). It runs only git; it runs none
of the departments' code and writes nothing in their repositories. Dates are committer
dates in America/New\_York; the author and committer dates differ in 17 of 1,750 commits. % ev: README.md "Environment and commands"
Re-running it at the pins gives byte-identical output.

\paragraph{Node lifecycles.}
For each repository with a board, every commit that changed \code{graph.yaml} is parsed
and diffed against its first parent. A node's time in a state is the first commit that
shows it there. Lead times are differences between those times. Node ids are replaced
by neutral ids, since some name the client or a song.

\paragraph{Merges and conflicts.}
A merge \emph{needed a hand resolution} if \code{git show --remerge-diff} is non-empty:
re-doing the merge automatically would give a different tree, so a person or session
changed something by hand. This is a stricter test than searching commit messages, which
mention a conflict in 43 studio merges. % ev: integrations.json repos.studio.merges_whose_message_mentions_a_conflict

\paragraph{Lead sessions.}
Only the studio names its sessions in commit subjects (``(lead-9)''). A commit is
attributed to a lead only if a parenthetical in its subject starts with exactly one lead
name; this attributes 221 of 1,355 studio commits (16\%). Concurrency is estimated by % ev: README.md "Lead sessions" (221, 16.3%)
treating a lead as active from its first to its last attributed commit. That span
includes idle time, and a lead whose commits never name it is invisible, so the estimate
can be wrong in both directions.

\paragraph{Lessons.}
The 61 dated lessons were coded once each by a Claude session acting as this % ev: incidents.json counts.dated_lessons
department's reproducer: kind (incident, practice, or the human's directive), one of
eleven cause classes, harm, fix, the rule produced, and the commit that added the lesson.
Copies of a lesson into another repository, and second lessons about the same event, are
linked to the original. A script checks that every lesson is coded, that dates and
sources match the repositories, and that no client or person is named. There was one
coder and no agreement measure (\S\ref{sec:threats}).

\paragraph{Committed versus observational.}
Most numbers come from committed history and can be re-mined. A few come from the
studio's \emph{uncommitted} state directory and deploy logs, read read-only with the
studio's agreement: PASS markers, number reservations and deploy logs. These files are
rewritten while the studio runs, so they are a snapshot taken at one moment, not history.
We mark them \obs{} in tables and say so in the text.

\paragraph{What was withheld.}
No client data, product names, measurements, media, audio, track names, session
transcripts, reviewer screenshots, commit subjects or e-mail addresses appear here or in
the evidence files. People are not named: the client's owner, the two executives and
other stakeholders are described by role.

% ---------------------------------------------------------------------------------------
\section{Results}\label{sec:results}

\subsection{Scale and timeline}

Figure~\ref{fig:timeline}a shows commits per day. Studio commits peaked at 552 on % ev: scale.json repos.studio.commits_per_day["2026-09-30"]
30 September. Of all 1,750 commits, 1,635 (93\%) carry a Claude co-author line; in the % ev: scale.json org.commits_with_claude_coauthor; derived.json org_coauthor_share
studio, 1,257 of 1,355 do, naming two models (744 and 513 commits). The other six % ev: scale.json repos.studio.commits_with_claude_coauthor, coauthor_commits_by_model
repositories were created on the last day, which alone holds 345 commits; in the hour % ev: scale.json org.commits_per_day["2026-10-02"]
from 05:00 on 2 October seven of the eight repositories committed. This is activity in % ev: derived.json max_repos_committing_in_one_hour
repositories, not a count of sessions: one session can commit to several repositories.

\begin{figure}[t]
  \centering
  \includegraphics[width=\linewidth]{timeline.pdf}
  \caption{(a) Commits per day in all eight repositories (committed history). (b) Studio
    lead sessions active at once, per hour. Solid: a session spans its first to last
    commit that names it. Dashed (\obs): a session starts when its lead number was
    reserved in the studio's state directory. Both assume a session is active for its
    whole span, idle time included.}
  \label{fig:timeline}
\end{figure}

The studio's sessions were numbered up to lead-13; 12 of them have attributed commits. % ev: leads.json repos.studio (13 names; 12 with attributed commits)
Across the organisation at least 22 lead sessions were named (13 in the studio, 9 in % ev: leads.json highest_number_by_department: 13+2+1+2+2+1+1 = 22
five other departments). Under the span assumption, at most five studio leads were active
at once (on 1 October), and at most four committed within the same clock hour % ev: leads.json repos.studio.concurrency_from_spans.max_overlap = 5; max_distinct_leads_committing_in_one_hour = 4
(Figure~\ref{fig:timeline}b). Starting each session when its number was reserved\obs{} % ev: studio_state.json concurrency_reserved_to_last_commit.max_overlap = 5
gives the same maximum of five. Spans range from 1.2 to 34 hours; the longest is the % ev: leads.json session_spans_from_attributed_commits (lead-8 1.163 h; lead-9 34.151 h)
session that held the operations duty.

\subsection{Node lifecycles}

Table~\ref{tab:lifecycle} and Figure~\ref{fig:lifecycle} give per-department lead times
from the boards. In the studio, 100 of 113 nodes reached \code{merged}; the median node % ev: lifecycle.json repos.studio.ever_merged, nodes
took 3.5 hours from first appearing as \code{todo} to merged (IQR 1.4--6.0, $n=62$) and % ev: lifecycle.json repos.studio.lead_time_hours.todo_to_merged_h (3.526; 1.377-6.042; n 62)
half an hour from claimed or building to merged ($n=71$). Twenty-two studio nodes were % ev: lifecycle.json ...claimed_to_merged_h (0.533, n 71); first_state_counts.building = 22
seeded already \code{building} when the board was created, so they have no \code{todo}
time. These are durations of board states, not effort: a node can wait in \code{todo}
for reasons that have nothing to do with the system, such as the human's priorities.

\begin{table}[t]
\centering
\small
\caption{Board lifecycles at the pins. Hours; median, with the interquartile range where
$n\ge 6$. ``todo$\to$merged'' starts at the first commit showing the node as \code{todo};
``claimed$\to$merged'' at the first showing it \code{claimed} or \code{building}.
product-in-picture and the gateway have no board.}
\label{tab:lifecycle}
\begin{tabular}{@{}lrrrll@{}}
\toprule
Repository & Nodes & Merged & Dropped & todo$\to$merged ($n$) & claimed$\to$merged ($n$) \\
\midrule
studio & 113 & 100 & 1 & 3.53 (1.38--6.04, 62) & 0.53 (71) \\ % ev: lifecycle.json repos.studio; README.md "Board lifecycles"
mir-triage (board branch) & 41 & 13 & 0 & 1.48 (0.05--1.74, 13) & 1.16 (12) \\ % ev: lifecycle.json repos.mir-triage
copy-in-product-picture & 37 & 13 & 0 & 1.44 (1.33--1.92, 6) & 1.05 (9) \\ % ev: lifecycle.json repos.copy-in-product-picture
product-in-video & 11 & 3 & 0 & 0.44 (3) & 0.08 (3) \\ % ev: lifecycle.json repos.product-in-video
research-to-publication & 3 & 2 & 0 & 0.17 (2) & 0.10 (2) \\ % ev: lifecycle.json repos.research-to-publication
design-manufacture-interface & 6 & 0 & 0 & -- & -- \\ % ev: lifecycle.json repos.design-manufacture-interface
\bottomrule
\end{tabular}
\end{table}

\begin{figure}[t]
  \centering
  \includegraphics[width=\linewidth]{lifecycle.pdf}
  \caption{Lead time of every merged node, per department, on a log scale: (a) first
    seen as \code{todo} to first \code{merged}; (b) first \code{claimed} or
    \code{building} to first \code{merged}. One dot per node; the bar is the median.
    Committed board history.}
  \label{fig:lifecycle}
\end{figure}

Two intervals need care. In the studio, \code{review} or \code{ready} to merged has a
median of 0.008 hours (about 30 seconds, $n=92$). The studio sets \code{ready} just % ev: lifecycle.json repos.studio.lead_time_hours.review_to_merged_h (0.008, n 92)
before running the integration script and reviews while a node is still
\code{building}, so this interval measures integration, not review. Merged to
\code{deployed} has a median of 0.087 hours (about five minutes, $n=69$). % ev: lifecycle.json repos.studio.lead_time_hours.merged_to_deployed_h (0.087, n 69)

The boards rarely move backwards. On the audio engine's trunk, four nodes went from % ev: lifecycle.json repos.mir-triage.state_regressions_on_trunk = 4
\code{review} back to \code{building} after an evaluator's FAIL, that is, a rework round;
the studio has one deliberate reopen and the other boards none. % ev: lifecycle.json repos.studio.state_regressions_on_trunk = 1

\subsection{Integration and merge conflicts}

Table~\ref{tab:merges} summarises merges. The studio made 130 branch-to-trunk merges and % ev: integrations.json repos.studio.merge_kinds["branch into trunk"]
208 merges of trunk into a branch: 53 by the integration script and 155 by hand. % ev: integrations.json repos.studio.trunk_into_branch_merges, trunk_into_branch_by_integrate_script; merge_kinds
Of all 351 studio merges, 51 (15\%) needed a hand resolution. Thirty-three of those % ev: integrations.json ...merges_needing_hand_resolution_remerge_diff; derived.json studio_hand_resolved_share
touched the board file, \code{graph.yaml}, and 30 touched nothing else. None was made by % ev: integrations.json hand_resolved_merges_touching_class; derived.json studio_hand_resolved_only_graph
the integration script, which stops on a conflict by design. Forty-seven were hand merges % ev: integrations.json merges_needing_hand_resolution_by_kind["trunk into branch (by hand)"] = 47
of trunk into a branch, the step a lead takes after the script stops: 30\% of the 155 % ev: derived.json studio_hand_resolved_share_of_hand_trunk_merges
such merges needed one.

\begin{table}[t]
\centering
\small
\caption{Merges at the pins. ``Hand-resolved'' means \code{git show --remerge-diff} is
non-empty. File classes overlap: one merge can touch several. Repositories with no
hand-resolved merge are omitted (product-in-picture, product-in-video,
research-to-publication, design-manufacture-interface, gateway).}
\label{tab:merges}
\begin{tabular}{@{}lrrrl@{}}
\toprule
Repository & Merges & Hand-resolved & Share & Files touched by the hand-resolved merges \\
\midrule
studio & 351 & 51 & 15\% & board 33, code 14, plans 9, generated 7, agent memory 6 \\ % ev: integrations.json repos.studio.hand_resolved_merges_touching_class
mir-triage & 37 & 8 & 22\% & code 6, plans 4 \\ % ev: integrations.json repos.mir-triage
copy-in-product-picture & 23 & 3 & 13\% & code 3, plans 2 \\ % ev: integrations.json repos.copy-in-product-picture
\bottomrule
\end{tabular}
\end{table}

The pattern is clear. Coordination through a shared board in git makes the board the
most conflicted file. The studio's lessons record the response: a node's state is set to
\code{built} on its branch, but \code{ready} and \code{merged} only on trunk, and setting % ev: incidents.json studio-14, studio-23
and committing a state became one locked command. Two other lessons show the cost of
resolving by hand: a conflict in a generated file was committed with its markers and
caught by the linter (2,258 errors), and conflict markers reached trunk once when a % ev: incidents.json studio-16 (as recorded in the lesson, not re-measured)
commit-all in the shared checkout swept up another lead's file (\S\ref{sec:incidents}). % ev: incidents.json studio-17
Code conflicts were fewer than board conflicts, which is what the \code{touches} rule is
meant to achieve; without a baseline we cannot say how many it prevented.

\subsection{The three-lead trial}\label{sec:trial}

The studio's plan records one deliberate trial, on the evening of 29 September (US % source: studio 6b04ffe, claude-plans/06-multi-lead.md "Trial record"
Eastern; 30 September UTC). Its record says three leads ran at once ``for about two
hours'' and lists ten nodes merged ``in parallel, no file, number or deploy collisions'',
three coordination conflicts each settled by the claim table or an agreed rule, and 0
CI minutes used because every gate was local.

We re-counted the trial from history. The window runs from the first lead's recorded
start to the merge just after the record commit: 1.26 hours. In it there were 111 % ev: integrations.json repos.studio.trial_2026_09_30.window_hours (1.262), commits (111)
commits, 29 merges, and 10 branch-to-trunk merges (the same ten the record lists), with % ev: trial_2026_09_30.merges (29), branch_into_trunk_merges (10)
9 board nodes reaching \code{merged}. Two merges needed a hand resolution, one on % ev: trial_2026_09_30.nodes_reaching_merged (9), merges_needing_hand_resolution (2)
\code{graph.yaml} and one on a shared plan. One board update was lost: a commit
overwrote a node's state and another set it ``ready again'' (\code{05ffd0f}, % ev: trial_2026_09_30.graph_state_reset_after_overwrite_commits; incidents.json hist-1
\code{94f8d2d}). Eight files other than shared indexes were changed by two or more of % ev: trial_2026_09_30.files_changed_by_more_than_one_merged_branch_excl_shared_indexes (8)
the merged branches (two code, two generated, four docs), and git merged all of them
cleanly. By the record's own times, three leads overlapped for only about a minute (the
second lead drained at about 01:40 UTC; the third started at about 01:39 UTC); for the rest
of the window there were two.

So the accurate summary is: \emph{no blocking conflicts, and no number or deploy
collisions}, rather than no collisions at all. The studio adopted this wording in a dated % source: studio 7870aca, claude-plans/06-multi-lead.md "Correction, 2026-10-02"
correction under the original record (\code{7870aca}). The ``0 CI minutes'' figure
cannot be checked from git; we quote it from the record and did not reproduce it.

\subsection{Review and deployment (observational)}\label{sec:review-deploy}

The review gate leaves little in git. Reviews happen while a node is \code{building}, and
their only lasting trace is a PASS marker in the studio's state directory. At the time of
the pinned commit there were 78 PASS markers for board nodes\obs{}. Each names a commit % ev: studio_state.json pass_markers.markers_for_board_nodes (78)
reachable from the pinned trunk, and each belongs to a node that was then \code{merged} % ev: studio_state.json pass_markers.per_node[].sha_reachable_from_main_pin
(22) or \code{deployed} (56). At the pin 104 studio nodes were merged or deployed. The % ev: studio_state.json pass_markers.node_state_at_pin; derived.json studio_nodes_merged_or_deployed_at_pin
difference is not a measure of skipped reviews: small nodes take a ``light'' gate, the
lead's own look at screenshots, which writes no marker. Review rounds per node cannot be
recovered in the studio. In the audio engine, the four rework rounds above are the only
recorded reviewer FAILs on a board.

The studio's live-history log\obs{} lists 116 deploys from 30 September to 2 October; in % ev: studio_state.json live_history.deploy_lines (116)
33 of them the deployed commit was not \code{main}, because built branches were overlaid % ev: studio_state.json live_history.deployed_sha_differs_from_main_sha (33)
(deploy at built). The auto-deployer's log\obs{} records 116 attempts: 111 deployed and % ev: studio_state.json autodeploy_log.line_kinds (deploying 116, deployed 111, deploy FAILED 5)
5 failed (two failed health checks, one database not at the migration head, two % ev: studio_state.json autodeploy_log.failed_deploys_by_reason
other), 13 migrations, and 8 waits for a quiet machine. Its one untimestamped line is % ev: studio_state.json autodeploy_log.line_kinds (migrated 13, deploy waits 8)
a shell syntax error in the deployer itself, which corroborates a lesson about rewriting
a running script in place (studio-15, \S\ref{sec:incidents}). % ev: studio_state.json autodeploy_log.quoted_line

\subsection{Kits}

Four departments that started on the last day from a kit have a board. Three of them % ev: derived.json hours_first_commit_to_first_merged_node
merged a first node within an hour of their repository's first commit (0.21, 0.64 and % ev: derived.json ...copy 0.21, piv 0.64, research 0.11
0.11 hours); the fourth had merged none by the pin. This shows that a kit makes the
machinery usable at once. It does not show how fast a department becomes productive:
first nodes were often the department's own setup, and some departments continued work
that began elsewhere.

% ---------------------------------------------------------------------------------------
\section{Incidents}\label{sec:incidents}

The eight playbooks hold 61 dated lessons: 25 in the studio, 10 in the audio engine, % ev: incidents.json counts.dated_lessons_by_repo
5 in still ads, 9 in product images, 5 in video, 3 in manufacturing, 2 in the gateway
and 2 here. Fifteen are copies of a lesson into another repository or a second lesson % ev: incidents.json counts.copied_or_same_event (15)
about the same event, which leaves 46 distinct: 33 incidents, 8 practices and 5 of the % ev: incidents.json counts.distinct_events_or_rules (46), distinct_by_kind
human's directives. Figure~\ref{fig:incidents} counts them by cause class, and
Table~\ref{tab:taxonomy} gives one example of each class with the rule it produced.
Appendix~\ref{app:incidents} lists every incident.

\begin{figure}[t]
  \centering
  \includegraphics[width=0.92\linewidth]{incidents.pdf}
  \caption{All 61 dated lessons by cause class: distinct incidents, distinct practices
    and directives, and copies (counted under the original's class). Hand-coded from the
    committed playbooks.}
  \label{fig:incidents}
\end{figure}

\begin{table}[t]
\centering
\small
\caption{Cause classes of the 33 distinct incidents, with one example each (scrubbed) and
the rule written back. Counts are distinct incidents; lesson ids refer to
Appendix~\ref{app:incidents}.}
\label{tab:taxonomy}
\begin{tabularx}{\linewidth}{@{}lrXX@{}}
\toprule
Cause class & $n$ & Example & Rule written back \\
\midrule
process/review & 10 & A ragged cut-out edge on a sample ad passed every review; the human and a stakeholder saw it (copy-3). & Gate cut-outs on an edge-quality metric. \\ % ev: incidents.json distinct_incidents_by_cause_class["process/review"]; copy-3
concurrency/coordination & 5 & A fresh lead found no claims, took every in-flight node and launched duplicate agents into the same worktrees (studio-8). & Claim before you dispatch. \\ % ev: ...["concurrency/coordination"]; studio-8
environment/OS & 5 & Worktrees under \code{/tmp} vanished on a reboot with an uncommitted fix round (studio-2). & Keep worktrees in the repo; commit work in progress. \\ % ev: ...["environment/OS"]; studio-2
deploy/live safety & 4 & A lead's edits, with an unfinished migration, landed in the main checkout; the auto-deployer would have migrated the live database (studio-20). & Name the worktree by absolute path in every command; the deployer reads only committed state. \\ % ev: ...["deploy/live safety"]; studio-20
resource contention & 3 & One browser-test run at default workers drove the load to 81 on 8 cores for 40 minutes (studio-22). & Check the load before any heavy run; cap workers. \\ % ev: ...["resource contention"]; studio-22 (numbers as recorded in the lesson)
tooling defect & 2 & The integration script counted trunk's own commits as unreviewed after merging trunk into the branch (studio-12). & Exclude commits reachable from trunk. \\ % ev: ...["tooling defect"]; studio-12
session lifecycle & 1 & Every lead died overnight at a lock screen, leaving claims and uncommitted files (studio-19). & Check liveness by process; split the work among fresh leads at once. \\ % ev: ...["session lifecycle"]; studio-19
communication & 1 & A lead's name relayed through two sessions was misread twice (pip-2). & Names and roles in the human's words only. \\ % ev: ...["communication"]; pip-2
content policy & 1 & Copying a reference ad's mechanic and naming its brand crossed the human's line (pip-4). & References are taste, not material. \\ % ev: ...["content policy"]; pip-4
privacy/data handling & 1 & A handoff recorded the client's product details in another repository's git (dmi-3). & Client numbers stay in the client's files; history purged. \\ % ev: ...["privacy/data handling"]; dmi-3
\bottomrule
\end{tabularx}
\end{table}

\paragraph{Process and review failures are the largest class.}
Ten distinct incidents are about planning, briefs, gates or reviews. In four incidents a % ev: incidents.json distinct_incidents_by_harm_class ("defect reached the user" 2 + "defect reached people" 2)
defect reached the human or a stakeholder. In two of them the defect had passed review % ev: copy-3, pip-5
(an edge judged at the wrong scale or on the wrong background); in one the page had no % ev: copy-5
review; and in one an end-to-end test exercised an action with a single input, and the % ev: studio-25
action failed on the human's real data in a demonstration. Reviewers that never edit are
independent, but they are not sufficient: what they look at, and on what data, matters
more than who looks. The rules produced are concrete: judge edges on the real
background with a number, and prove an action by doing it the way the user will.

\paragraph{One machine is a shared resource, and it fails as one.}
Resource contention, environment and deploy safety together account for 12 of the 33 % ev: 3 + 5 + 4 from distinct_incidents_by_cause_class
incidents. These are failures that a single coding session on its own machine does not
meet. The most-copied lesson in the organisation is one of them: a default-parallel test
run starved the laptop while the live app ran on it, and a regex test filter matched the
worktree's own name and ran the whole suite (studio-22). It was copied into five other % ev: incidents.json counts.times_each_original_was_copied["studio-22"] = 5
repositories. Later, three departments on the laptop at once drove the load to 29--44 % ev: incidents.json mir-8 (as recorded in the lesson)
(mir-8); the response was an organisation-wide cap: at most about six busy threads per
department, nothing heavy while the one-minute load is above 15, and long jobs under a
shared lock.

\paragraph{Sessions die, and leave state behind.}
Sessions end without warning: at a lock screen, at a reboot, or near the end of their
context. Claims, locks and uncommitted files outlive them. The rules are to commit work
in progress, keep worktrees inside the repository, check liveness by process rather than
by claim, and stop taking work at about 45\% of the context and write a handoff (a % ev: incidents.json studio-11 (practice)
practice, studio-11).

\paragraph{Lessons spread.}
A quarter of all lessons (15 of 61) are copies. Seven originals were copied: three from % ev: derived.json lessons_copied_share (0.246); incidents.json counts.times_each_original_was_copied (7 keys)
the studio, two from still ads, and one each from the audio engine and manufacturing. % ev: times_each_original_was_copied: studio-22, studio-20, studio-25, copy-3, copy-5, mir-8, dmi-3
The copies travel with the kits and with the organisation-wide rules file that every
session loads.

% ---------------------------------------------------------------------------------------
\section{Discussion}\label{sec:discussion}

\paragraph{What worked, as far as the record shows.}
Several leads shared one repository for days, and none of the 53 merges the integration % ev: integrations.json repos.studio.trunk_into_branch_by_integrate_script (53)
script made needed a hand resolution. Number and deploy collisions do not appear in the
history or the lessons. Every incident we coded ends in a rule written back into a
repository. None of this shows that the design is better than an alternative; it shows
that it held together under real use.

\paragraph{Lessons for practitioners.}
\begin{itemize}
  \item \emph{Expect the board to be the hot spot.} If coordination state lives in a
    file in git, that file will be the most conflicted file. Change it only on trunk,
    only under the lock, and set and commit it in one step.
  \item \emph{Schedule by declared file sets.} Declaring \code{touches} globs and
    refusing to offer overlapping nodes is cheap and kept code conflicts below board
    conflicts here.
  \item \emph{Make the gate stop, not resolve.} The integration script never resolved a
    conflict. Hand resolutions happened on branches, where a mistake is caught before
    trunk, although twice one was not (studio-16, studio-17).
  \item \emph{Budget the machine.} With several sessions on one laptop, CPU, the test
    lock and the live app are shared. A load cap, a heavy-job lock and ``check the load
    first'' were the most-copied rules.
  \item \emph{Plan for sessions to die.} Commit work in progress; keep claims and
    worktrees discoverable; hand off before the context runs out.
  \item \emph{Review what the user will see, the way the user will see it.} Independent
    reviewers missed defects when they judged at the wrong scale, on the wrong
    background, or with toy inputs.
  \item \emph{Write the lesson where the next session will read it.} A dated rule in the
    playbook reaches every future session of that department, and kits copy it to new
    ones.
\end{itemize}

\paragraph{Governance: direction is not approval.}
The split between direction and approval was tested once on record. When the human's
grant of autonomy was relayed from another session, leads refused to act on it until the % ev: incidents.json gateway-1 (practice)
human confirmed it in their own sessions, which was the correct behaviour (gateway-1).
The human later gave approvals directly, in each lead's session. The split is simple: a
relay may change \emph{what} a lead works on, never \emph{what it is allowed to do}. It
stops a chain of sessions from laundering a permission that no single session holds.
This is close to the idea of autonomy as a design decision separate from
capability~\citep{feng2025levels}; here it is applied to delegation between agents.

\paragraph{The record was corrected.}
The trial record's ``no collisions'' was stronger than the history supports. After our
re-count, the studio added a dated correction under its original sentence rather than
rewriting it (\code{7870aca}). We take that as a small sign that committed, dated records
can be audited and fixed in place, which is what made this study possible at all.

\paragraph{What may generalise.}
The design uses nothing specific to one vendor's agents: git, directories as locks, a
YAML board and shell scripts. The failure classes we saw (shared-machine contention,
dying sessions, coordination state in git) should appear in any setup where several
agents share one machine and one repository. What may not generalise is the scale (one
human, one laptop, under a week), the organisation's taste for speed over caution (for
example deploying unreviewed work to the human's phone), and the products.

% ---------------------------------------------------------------------------------------
\section{Related work}\label{sec:related}

\paragraph{Multi-agent LLM systems for software.}
MetaGPT~\citep{hong2023metagpt} and ChatDev~\citep{qian2024chatdev} organise agents into
roles that follow a software process and evaluate them on generated tasks. He et
al.~\citep{he2024llmmas} survey LLM-based multi-agent systems for software engineering.
SWE-agent~\citep{yang2024sweagent} studies the interface a single agent needs to work in a
repository. Our system uses roles too (leads, builders, reviewers), but the unit of study
is a live organisation over days, not a run on a task.

\paragraph{Coordination and failures of agent teams.}
Destefanis and Aste~\citep{destefanis2026coordinate} measure coordination in 1,902 runs of
multi-agent coding teams, varying team size, structure and file policy. Cemri et
al.~\citep{cemri2025mast} build a taxonomy of 14 failure modes from annotated traces of
multi-agent systems. De Oliveira et al.~\citep{oliveira2026masexperience} report on
building with open-source multi-agent frameworks for a software engineering task. EvoGit~\citep{huang2025evogit}
uses git's history as the coordination substrate for decentralised agents. Our incidents
differ in kind from failures inside a trace: most are about the shared machine, sessions
ending, deployment and the human's gates. Our data are observational and come from one
organisation, so they complement these controlled studies rather than test them.

\paragraph{Practice.}
Practitioners describe running many agents in parallel worktrees with a human
orchestrating~\citep{mason2026agents}. We add a measured account of one such setup,
including its conflicts and incidents.

\paragraph{Merge conflicts and their avoidance.}
Studies of human teams find that merge conflicts are common and costly: Brindescu et
al.~\citep{brindescu2020merge} link conflicts to later defects, and Nelson et
al.~\citep{nelson2019lifecycle} describe how developers plan for, resolve and avoid them.
Proactive approaches detect conflicts early~\citep{brun2011crystal} or schedule tasks to
avoid them~\citep{kasi2013cassandra}. The \code{touches} rule is a crude form of the
latter. Our finding that the coordination file itself is the main source of conflicts
is specific to putting the coordination state in the same repository.

\paragraph{Incident analysis and measurement.}
The lessons resemble blameless postmortems, where each incident produces an action that
prevents a repeat~\citep{lunney2016postmortem}. Our lessons are much shorter than
postmortems and written by the sessions themselves. One department measured how much a
shared, busy laptop moves benchmark results; the method follows Kalibera and
Jones~\citep{kalibera2013benchmarking} and is not part of this paper's claims.

% ---------------------------------------------------------------------------------------
\section{Threats to validity}\label{sec:threats}

\paragraph{No speed, productivity or cost claims.}
Nothing recorded the tokens, wall-clock time or money each session used, and there is no
baseline of one lead, or of a human team, doing comparable work. We therefore claim
nothing about speed-up, productivity or cost. Commit counts measure activity, not output.

\paragraph{Construct validity.}
Board states are what sessions wrote, not what happened; a state can be set late or
skipped, and the first commit showing a state is our proxy for when it was reached.
``Hand resolution'' via \code{--remerge-diff} also counts deliberate edits made during a
merge. The review$\to$merged interval measures integration, not review. Concurrency rests
on the first-to-last-commit span: idle time inflates it, and leads that never named
themselves in a commit subject are invisible; outside the studio it cannot be computed.

\paragraph{The lessons.}
Lessons are self-reported by the sessions that had the incidents. They are not a complete
incident log: an incident nobody wrote down is missing, and the count depends on how
eagerly each department writes lessons. They were coded by one coder, a Claude session,
with no second coder and no agreement measure; the class boundaries are ours. Lesson dates
are as written and follow UTC, so some differ by a day from the commit that added them.

\paragraph{Observational data.}
PASS markers, reservations and deploy logs\obs{} come from the studio's uncommitted state,
read once. They are a snapshot that the studio rewrites as it runs, and cannot be
re-mined later.

\paragraph{The trial.}
The trial was short (1.26 hours), had three leads for only about a minute, and was run by % ev: integrations.json trial_2026_09_30.window_hours
the system's own sessions. Its record overstated the result until it was corrected
(\code{7870aca}). ``0 CI minutes'' is quoted, not reproduced.

\paragraph{External validity and the observer.}
One human, one laptop, one model family and under a week. The organisation changed
during the period: six of eight repositories are less than a day old at the pins, and % ev: derived.json repos_first_commit_on_2026_10_02
rules were added as incidents happened. The author owns and steered the system, and
Claude, which built and runs most of it, also helped mine and write up the evidence. We
reduced this risk by deriving every number with committed scripts at pinned commits,
having the source department check and correct its own record, and stating what we could
not measure; it is not removed.

% ---------------------------------------------------------------------------------------
\section{Conclusion}\label{sec:conclusion}

One engineer ran eight departments of concurrent Claude lead sessions on one laptop for
five and a half days, coordinated through git and a handful of files. The committed
history shows real concurrency in the studio (up to five leads under our span
assumption), a board file that absorbed most of the conflicts, an integration gate that
never had to resolve one, and a steady stream of concrete incidents, each turned into a
rule that the next session reads. The most useful rules were about the shared machine,
dying sessions and where coordination state lives, not about the agents' code. Whether
the design is faster or cheaper than the alternatives is an open question this case
cannot answer; a controlled comparison with recorded cost is the obvious next study.

% ---------------------------------------------------------------------------------------
\section*{AI assistance}
The system studied here is built and operated largely by Claude sessions (Anthropic).
Claude also assisted with this paper: mining the repositories, coding the lessons,
analysing the evidence and drafting the text, under the author's direction. The author
chose what to study, reviewed the analysis and is responsible for the content.

% TODO(user): acknowledgements, if any, are the author's call and not yet made.
% TODO(user): data availability. The evidence files and scripts accompany this draft; the
% department repositories are private. Decide what is released with a preprint.

\bibliographystyle{plainnat}
\bibliography{refs}

% ---------------------------------------------------------------------------------------
\appendix
\section{Reproduction}\label{app:repro}

Everything in this paper builds from sources in \code{papers/org-of-one/}:
\begin{itemize}
  \item \code{make clean all} regenerates the four figures from their scripts
    (\code{figures/*.py}, Python 3.11 under \code{uv}, wheels only), the appendix table
    from \code{tables/incidents\_appendix.py}, and the PDF (TeX Live 2022, \code{latexmk}).
  \item The evidence is re-mined with
    \code{nice -n 10 uv run --python 3.11 --script mine.py --clones \$CLONES --pins sources.json}
    in \code{evidence/}, where \code{\$CLONES} is any scratch directory. It makes bare
    clones and runs only git. The observational snapshot comes from
    \code{studio\_state.py}; \code{check\_incidents.py} verifies the coded lessons;
    \code{derived.py} computes the few shares and spans stated here from the other files.
  \item Pins (\code{evidence/sources.json}, 2026-10-02 06:40 US Eastern): studio
    \code{af5aaf2}, mir-triage \code{2599ceb} (board branch \code{f0dacfb}),
    copy-in-product-picture \code{48f062f}, product-in-picture \code{c0b095d},
    product-in-video \code{8b6f38b}, design-manufacture-interface \code{8cc1238},
    gateway \code{a9a58bd}, research-to-publication \code{887dd2c}. The studio's trial
    correction is read at \code{7870aca}, which descends from the studio pin.
\end{itemize}

\section{Incident list}\label{app:incidents}
\input{build/tables/incidents_appendix.tex}

\end{document}
